\documentclass{article}
\usepackage[margin=1in]{geometry}
\usepackage{longtable,graphicx}
\begin{document}
\section*{Five-flow source recovery evaluation}
PASS means no reproducible final divergence was detected within the recorded valid inputs and budget; it is not a universal equivalence proof.
\begin{longtable}{llllllllllllllll}
Flow & N & Eligible N & First-pass RSR & Final RSR@R & Program Behavioral Pass Rate & Initial Behavioral Pass Rate & Final Behavioral Pass Rate & Input Match Macro & Input Match Micro & Counterexample Detection Rate & Re-executability Rate & Compilation Repair Gain & Behavioral Repair Gain & Mean Tokens & Mean Runtime \\ \hline
B1 & 40 & 40 & 0/1 (0.0\%) & 0/1 (0.0\%) & N/A & N/A & N/A & N/A & N/A & N/A & 0/40 (0.0\%) & N/A & N/A & 34603 & 15.1s \\
B2 & 40 & 40 & 13/15 (86.7\%) & 13/15 (86.7\%) & 3/13 (23.1\%) & 3/13 (23.1\%) & 3/13 (23.1\%) & 27.6\% & 29.9\% & 10/13 (76.9\%) & 13/40 (32.5\%) & N/A & N/A & 63591 & 29.5s \\
F1 & 40 & 40 & 4/4 (100.0\%) & 4/4 (100.0\%) & 4/4 (100.0\%) & 4/4 (100.0\%) & 4/4 (100.0\%) & 100.0\% & 100.0\% & 0/4 (0.0\%) & 4/40 (10.0\%) & +0.0 pp & +0.0 pp & 1148440 & 78.1s \\
F2 & 40 & 40 & N/A & N/A & N/A & N/A & N/A & N/A & N/A & N/A & 0/40 (0.0\%) & N/A & N/A & 231926 & 16.4s \\
F3 & 40 & 40 & 6/6 (100.0\%) & 6/6 (100.0\%) & 3/6 (50.0\%) & 3/6 (50.0\%) & 3/6 (50.0\%) & 62.6\% & 76.6\% & 3/6 (50.0\%) & 6/40 (15.0\%) & +0.0 pp & +0.0 pp & 755006 & 81.2s \\
\end{longtable}
\section*{Source quality: accepted Recovered C Source only}
Readability is an independent 1-to-5 assessment of Variables, Loops, Conditions, Logic flow, and Structural integrity. It is not a correctness metric.
\begin{longtable}{llllllllll}
Flow & Accepted & Evaluated & Coverage & Variables & Loops & Conditions & Logic flow & Structural integrity & Overall \\ \hline
B1 & 0 & 0 & 0/0 (N/A) & N/A & N/A & N/A & N/A & N/A & N/A \\
B2 & 3 & 3 & 3/3 (100.0\%) & 3.67 & 4.67 & 4.67 & 4.67 & 4.67 & 4.47 \\
F1 & 4 & 4 & 4/4 (100.0\%) & 4.00 & 5.00 & 4.50 & 5.00 & 5.00 & 4.70 \\
F2 & 0 & 0 & 0/0 (N/A) & N/A & N/A & N/A & N/A & N/A & N/A \\
F3 & 3 & 3 & 3/3 (100.0\%) & 4.00 & 5.00 & 5.00 & 5.00 & 4.67 & 4.73 \\
\end{longtable}
\begin{figure}[p]\centering\includegraphics[width=.9\linewidth]{figures/overall_performance.pdf}\caption{Primary recovery outcomes with executable availability over all eligible samples.}\end{figure}
\begin{figure}[p]\centering\includegraphics[width=.9\linewidth]{figures/compilation_performance.pdf}\caption{Compilation success before and after repair.}\end{figure}
\begin{figure}[p]\centering\includegraphics[width=.9\linewidth]{figures/behavioral_performance.pdf}\caption{Behavioral outcomes over conclusive campaigns.}\end{figure}
\begin{figure}[p]\centering\includegraphics[width=.9\linewidth]{figures/repair_effectiveness.pdf}\caption{Behavioral repair effectiveness; B1, B2, F2 repair fields are N/A.}\end{figure}
\begin{figure}[p]\centering\includegraphics[width=.9\linewidth]{figures/iterative_feedback_vs_one_shot.pdf}\caption{Matched-evidence re-executability comparison of iterative feedback against one-shot reconstruction.}\end{figure}
\begin{figure}[p]\centering\includegraphics[width=.9\linewidth]{figures/cumulative_compile_success_by_round.pdf}\caption{Cumulative executable creation by compile attempt.}\end{figure}
\begin{figure}[p]\centering\includegraphics[width=.9\linewidth]{figures/cumulative_behavioral_success_by_round.pdf}\caption{Cumulative behavioral pass by campaign.}\end{figure}
\begin{figure}[p]\centering\includegraphics[width=.9\linewidth]{figures/final_status_breakdown.pdf}\caption{Final status distribution without collapsing INCONCLUSIVE into failure.}\end{figure}
\begin{figure}[p]\centering\includegraphics[width=.9\linewidth]{figures/failure_taxonomy.pdf}\caption{Compilation, behavioral, and inconclusive taxonomies.}\end{figure}
\begin{figure}[p]\centering\includegraphics[width=.9\linewidth]{figures/sample_flow_status_heatmap.pdf}\caption{0=fail, 1=unresolved, 2=pass by paired sample and flow.}\end{figure}
\begin{figure}[p]\centering\includegraphics[width=.9\linewidth]{figures/sample_flow_match_heatmap.pdf}\caption{Input match (\%) by paired sample and flow.}\end{figure}
\begin{figure}[p]\centering\includegraphics[width=.9\linewidth]{figures/ablation_forest_plot.pdf}\caption{Paired re-executability effects with bootstrap 95\% CI.}\end{figure}
\begin{figure}[p]\centering\includegraphics[width=.9\linewidth]{figures/ablation_win_tie_loss.pdf}\caption{Paired win/tie/loss counts for re-executability.}\end{figure}
\begin{figure}[p]\centering\includegraphics[width=.9\linewidth]{figures/tokens_vs_e2e.pdf}\caption{Total tokens versus Re-executability outcome.}\end{figure}
\begin{figure}[p]\centering\includegraphics[width=.9\linewidth]{figures/runtime_vs_behavioral_pass.pdf}\caption{Runtime (seconds) versus Behavioral pass outcome.}\end{figure}
\begin{figure}[p]\centering\includegraphics[width=.9\linewidth]{figures/llvm_reduction_summary.pdf}\caption{LLVM simplification metrics; these are not correctness claims.}\end{figure}
\begin{figure}[p]\centering\includegraphics[width=.9\linewidth]{figures/llvm_reduction_vs_behavior.pdf}\caption{Association view only; LLVM reduction is not a semantic oracle.}\end{figure}
\begin{figure}[p]\centering\includegraphics[width=.9\linewidth]{figures/stage_completion_funnel.pdf}\caption{Stage completion rates with explicit numerators and denominators in the report.}\end{figure}
\end{document}
